\section{Discussion}
We propose ReMIA, a new privacy metric inspired by MIAs but based on a new attack model that requires significantly less computational and data resources than state-of-the-art MIAs, making it a practical tool for evaluating the privacy of synthetic data generators in realistic settings.
Across multiple datasets and generators, our experiments show that ReMIA correlates with state-of-the-art MIAs and achieves competitive sensitivity relative to lower-cost alternatives to MIAs based on shadow modeling.
Moreover, we find empirical evidence that synthetic data generated by SDGs can obtain a better privacy-quality trade-off than noise-based anonymized data according to different privacy metrics.


\paragraph{Limitations.}
Like most privacy metrics, ReMIA depends on implementation choices, including classifier design and training protocol.
Moreover, extending comparisons to a broader range of SDGs, anonymization methods, datasets, training sizes, and attack settings remains limited by the cost of running smMIAs and the cost of training powerful SDGs. In particular, a more extensive exploration of the effect of training size on privacy risk is an interesting direction for future work.

\paragraph{Future work.}
Future work should focus on studying the theoretical connections between ReMIA, overfitting, and MIAs, to better understand when ReMIA or detection of overfitting are good proxies for privacy risk.
On the methodological side, stronger attacker models may increase sensitivity while preserving computational efficiency.
Finally, the use of outliers as in \citet{meeus2023achilles} for improving the performance of ReMIA is an interesting direction to explore.

\paragraph{Broader impact.}
As a practical tool for evaluating the privacy of synthetic data generators, ReMIA can help researchers and practitioners to better detect and mitigate privacy risks associated with synthetic data. By providing a more computationally accessible metric for privacy evaluation, it may encourage the development and adoption of safer synthetic data generation techniques, ultimately contributing to the responsible use of synthetic data in various applications. On the other hand, like any empirical metric, it may not capture all aspects of privacy risk, and its misuse or misinterpretation could lead to a false sense of security.
